% Part: sets-functions-relations
% Chapter: functions
% Section: kinds

\documentclass[../../../include/open-logic-section]{subfiles}

\begin{document}

\olfileid{sfr}{fun}{kin}
\olsection{فنکشناں دیاں قسماں (Kinds of Functions)}

\begin{explain}
ساڈے کم آوے گا جے اسی عام طور تے ملن والیاں فنکشناں (functions) دیاں کجھ
قسماں دی درجہ بندی متعارف کرا دیئے۔

شروع وچ اسی اوہناں فنکشناں اُتے غور کر سکدے آں جنہاں دے ہدف سیٹ
(codomain) دا ہر رکن فنکشن دی اک قدر ہوندا اے۔ ایہو جہے فنکشناں نوں
!!{surjective} (ہر ہدف قدر لین والے؛ surjective) آکھیا جاندا اے، تے
\olref{fig:surjective} وچ دتے خاکے وانگ وکھایا جا سکدا اے۔

\begin{figure}
  \olasset{assets/diagrams/surjective.tikz}
  \caption{!!^a{surjective} فنکشن دے ہدف سیٹ (codomain) دا ہر !!{element}
    (element/member) اوہدی
    اک قدر ہوندا اے۔}
  \ollabel{fig:surjective}
\end{figure}
\end{explain}

\begin{defn}[!!^{surjective} فنکشن (ہر ہدف قدر لین والا فنکشن)]
اک فنکشن $f \colon A \rightarrow B$ \emph{!!{surjective}} عین اوس
ویلے اے جدوں $B$ وی~$f$ دا حاصل قدراں دا سیٹ (range) ہووے، یعنی ہر
$y \in B$ لئی گھٹ توں گھٹ اک $x \in A$ ہووے جیہدے لئی~$f(x) = y$؛
علامتاں وچ:
\[
  (\forall y \in B)(\exists x \in A)f(x) = y.
\]
اسی ایہو جہے فنکشن نوں $A$ توں $B$ ول !!a{surjection} (سرجیکشن)
آکھدے آں۔
\end{defn}

\begin{explain}
جے تسی وکھاؤنا چاہندے او کہ $f$ اک !!a{surjection} اے، تاں ایہ
وکھاؤنا پئے گا کہ $f$ دے ہدف سیٹ (codomain) دی ہر شے، $f(x)$ دی قدر
(output/value) اے، کسے انپٹ (input) $x$ لئی۔

خیال رکھو کہ ہر فنکشن اک !!a{surjection} (سرجیکشن) \emph{پیدا کردا} اے۔ آخر،
دتے ہوئے فنکشن $f \colon A \to B$ لئی $f' \colon A \to \ran{f}$ نوں
$f'(x) = f(x)$ راہیں تعریف کرو۔ کیوں جے $\ran{f}$ نوں
\emph{تعریف ای} $\Setabs{f(x) \in B}{x \in A}$ کیتا گیا اے، ایس
لئی فنکشن $f'$ دا !!a{surjection} (سرجیکشن) ہونا یقینی اے۔
\end{explain}

\begin{explain}
ہُن ہر فنکشن ہر ممکن انپٹ (input) نوں اک اکلوتے آؤٹ پٹ (output) نال جوڑدا اے۔ پر
ایہو جہے فنکشن وی نیں جیہڑے دو وکھ وکھ انپٹاں نوں کدے اکو آؤٹ پٹ
نال نہیں جوڑدے۔ ایہو جہے فنکشناں نوں !!{injective} (اک بہ اک؛ injective)
آکھیا جاندا اے،
تے \olref{fig:injective} وچ دتے خاکے وانگ وکھایا جا سکدا اے۔
\begin{figure}
  \olasset{assets/diagrams/injective.tikz}
  \caption{!!^a{injective} فنکشن (اک بہ اک فنکشن) دو وکھ وکھ دلیلاں نوں کدے اکو قدر
    نال نہیں جوڑدا۔}
  \ollabel{fig:injective}
\end{figure}
\end{explain}

\begin{defn}[!!^{injective} فنکشن (اک بہ اک فنکشن)]
اک فنکشن $f \colon A \rightarrow B$ \emph{!!{injective}} عین اوس
ویلے اے جدوں ہر $y \in B$ لئی ودھ توں ودھ اک $x \in A$ ہووے جیہدے
لئی~$f(x) = y$۔ اسی ایہو جہے فنکشن نوں $A$ توں~$B$ ول
!!a{injection} (انجیکشن) آکھدے آں۔
\end{defn}

\begin{explain}
جے تسی وکھاؤنا چاہندے او کہ $f$ اک !!a{injection} (انجیکشن) اے، تاں ایہ
وکھاؤنا پئے گا کہ کنہاں وی !!{element}s (elements) $x$ تے $y$ لئی، جیہڑے
$f$ دے دائرۂ تعریف (domain) وچ نیں، جے $f(x)=f(y)$ ہووے، تاں $x=y$ ہووے۔
\end{explain}

\begin{ex}
مستقل فنکشن (constant function) $f\colon \Nat \to \Nat$، جیہڑا $f(x) = 1$ نال دتا
گیا اے، نہ !!{injective} (injective) اے، نہ !!{surjective} (surjective)۔

شناختی فنکشن (identity function) $f\colon \Nat \to \Nat$، جیہڑا $f(x) = x$ نال دتا
گیا اے، !!{injective} وی اے تے !!{surjective} وی۔

جانشین فنکشن (successor function) $f \colon \Nat \to \Nat$، جیہڑا $f(x) = x+1$ نال
دتا گیا اے، !!{injective} اے پر !!{surjective} نہیں۔

ایویں صورتاں دے حساب نال تعریف کیتا فنکشن (function defined by cases)
$f \colon \Nat \to \Nat$:
\[
  f(x) =
  \begin{cases}
    \frac{x}{2} & \text{جے $x$ جفت اے} \\
    \frac{x+1}{2} & \text{جے $x$ طاق اے۔}
  \end{cases}
\]
!!{surjective} (surjective) اے، پر !!{injective} (injective) نہیں۔
\end{ex}

\begin{explain}
اکثر سانوں اوہناں فنکشناں اُتے غور کرنا ہوندا اے جیہڑے
!!{injective} وی ہون تے !!{surjective} وی۔ اسی ایہو جہے فنکشناں نوں
!!{bijective} (اک بہ اک تے ہر ہدف قدر لین والے؛ bijective) آکھدے آں۔ اوہ
\olref{fig:bijective} وچ وکھائے فنکشن
وانگ ہوندے نیں۔ !!^{bijection}s نوں کدی کدی
\emph{اک بہ اک مطابقتاں (one-to-one correspondences)} وی آکھیا جاندا اے، کیوں جے اوہ ہدف سیٹ (codomain) دے
عنصراں نوں دائرۂ تعریف دے عنصراں نال اکلوتے طور تے جوڑدے نیں۔
\begin{figure}
  \olasset{assets/diagrams/bijective.tikz}
  \caption{!!^a{bijective} فنکشن (bijective function) ہدف سیٹ دے عنصراں نوں دائرۂ تعریف
    دے عنصراں نال اکلوتے طور تے جوڑدا اے۔}
  \ollabel{fig:bijective}
\end{figure}
\end{explain}

\begin{defn}[!!^{bijection} (بائجیکشن)]
اک فنکشن $f \colon A \to B$ \emph{!!{bijective}} عین اوس ویلے اے
جدوں اوہ !!{surjective} وی ہووے تے !!{injective} وی۔ اسی ایہو جہے
فنکشن نوں $A$ توں~$B$ ول (یا $A$ تے~$B$ دے وچکار) !!a{bijection}
(بائجیکشن) آکھدے آں۔
\end{defn}

\end{document}
